\documentclass{ximera}

\title{Exponential Function Derivatives Activity Facilitator Guide}
\author{MATH 425: Calculus I}

\begin{document}
\begin{abstract}
    Working with the peers in your group, solve the following problems. Make sure to show and justify all your work. Make sure everyone in the group understands the solution and participates. Be prepared to report your answers to the whole class. 
\end{abstract}
\maketitle

\textbf{Student Learning Objectives}\\
Students will be able to:
    \begin{itemize}
        \item Find the derivative of functions of the form $f(x)=e^x$ using the constant multiple, sum, and difference rules  (problems 2-6).
        \item Find the derivative of functions of the form $f(x)=e^{g(x)}$ using the product, quotient, and chain rules (problems 7-8).
        \item Interpret the derivative of an exponential function as the slope of the tangent line, or in a physical context.
        \item Write the equation of tangent line to a simple function involving an exponential function.
        \item Hypothesize regarding the connection between $e^tln(b)$ and $b^t$.
    \end{itemize}

\textbf{Prerequisite Knowledge:} This activity, as a whole, expects that students know the derivative of the exponential function $f(x)=e^x$, and the derivative rules.  Students are also expected to know the connection of a derivative to a tangent line. \\
This activity may be used in its entirety or in parts (solo or appended to other activities) as follows: \\
\begin{itemize}
    \item Problem 1 is a review/exploration of the concept of the exponential function and its connection to logarithms, although logarithms are not covered in this activity.  
    \item Problems 2-6 can be completed after only the basic derivative rules: power, constant multiple, sum and difference. 
    \item Problems 7-8 include the product, quotient and chain rules.

\begin{exercise}
    Open this Desmos worksheet: \\
    \begin{center}
        \desmos{hpym0hoois}{800}{600}
    \end{center}

    \begin{enumerate}
        \item By experimenting with the value of $k$, find a value of $k$ so that the graph of $p(t)=f(kt)=e^{kt}$ appears to align with the graph of $g(t)=3^t$.  What is the value of $k$?\\
        \textcolor{blue}{ $k$ looks to be approximately 1.1.}

        \item Similarly, experiment to find a value of $k$ so that the graph of $p(t)=f(kt)=e^{kt}$ appears to align with the graph of $h(t)=(\frac13)^t$.  What is the value of $k$?\\
        \textcolor{blue}{ This $k$ looks to be about -1.1}

        \item For the value of $k$ you determined in (a), compute $e^k$.  What do you observe? \\
        \textcolor{blue}{ $e^{1.1}\approx 3.004$, which is close to 3.  }
        \textcolor{orange}{Students should include something here stating that $e^k$ is about 3. }
   
        \item For the value of $k$ you determined in (a), compute $e^k$.  What do you observe?  \\
        \textcolor{blue}{ $e^{-1.1}\approx 0.339\approx \frac13$.  Note here that $e^k\approx \frac{1}{3}$}
   
        \item Given any exponential function of the form $b^t$, do you think it's possible to find a value of $k$ so that $p(t)=f(kt)=e^{kt}$ is the same function as $b^t$?  Why or why not?\\
        \textcolor{blue}{ We can find this number $k$ such that $e^{kt}=b^t$ when $b>0$, since we just need that $e^k=b$, and the range of an exponential function is all positive real numbers.}\\
        \textcolor{orange}{Students will likely need help putting their tought into mathematical terms here.  Keep digging and asking ``why?'' to make sure they are understanding the deeper concepts here.}
   
        \item  We'll be talking about logarithms and their connection to exponential functions soon. 
        It happens to be true that $ln(3)\approx 1.1$ and $\ln(\frac{1}{3})\approx -1.1$. With this information, can you set up an equation to find such a $k$ so that $p(t)=f(kt)=e^{kt}$ is the same function as $b^t$? \\
        \textcolor{blue}{ If $e^k=b$, then $\ln(b)=k$.}\\
        \textcolor{orange}{This is simply asking students to create for themselves the connection that $e^{\ln(b)}=b$.  Although we have not yet talked about logarithms, this rule will form an important part of later topics.   }
   
   
    \end{enumerate}  
   
    Problem adapted from Boelkins, et al., \textit{Active Prelude to Calculus}.
\end{exercise}

\begin{exercise}
    Calculate the derivative for the following functions using the derivative rules we've discussed in class so far: 
        \begin{enumerate}
                \item $g(x)=3e^x-\frac{4}{\sqrt{x}}$\\
                \textcolor{blue}{ $g'(x)=3e^x+2 x^{-\frac32}$}
    
                \item $f(x)=\sqrt[6]{x^7} -\frac{e^x}{3}$\\
                \textcolor{blue}{ $f'(x)=\frac{7}{6}x^{\frac16}-\frac13e^x$}
    
                \item $h(t)=4et^2$\\
                \textcolor{blue}{ $h'(t)=8et$}
    
        \end{enumerate}
\end{exercise}

\begin{exercise}
    Find the equation of the line tangent to the graphs of $f(x)=e^x$ at $x=0$.  Use Desmos to graph the function and your tangent line and check your work.\\
    \textcolor{blue}{ $f'(x)=e^x$ so $f'(0)=e^0=1$\\
    So the tangent line is $y-1=1(x-0)$ or $y=x+1$}
    \begin{center}
        \desmos{dwxditai89}{800}{600}
    \end{center}
    
\end{exercise}

\begin{exercise}
    Suppose $\displaystyle \lim_{x\rightarrow 2} f(x)=-1$. 
    What is $\displaystyle \lim_{x\rightarrow 2}e^{f(x)}$? (Hint: how do continuous functions and limits interact?)\\
    \textcolor{blue}{ $\lim_{x\rightarrow 2} e^{f(x)}=e^{\lim_{x\rightarrow 2} f(x)}=e^{-1}=\frac1e$.}\\
    \textcolor{orange}{This problem is not meant to be a trick.  This is just a reminder for students how continuous functions and limits can interact, now that we know some non-algebraic continuous functions and we will use this fact somewhat regularly.}
\end{exercise}


\begin{exercise}
    At what point on the curve $y=1-e^x$ is the tangent line parallel to $-x-y=5$?   Find the equation of the tangent line at this point.  Check your graph by graphing the curve and both lines on Desmos.\\
     \textcolor{blue}{$y'=-e^x$ and the line can be written as $y=-x-5$, which has slope -1.  \\
    Then we know that we want to know when $y'=-1$, or $-e^x=-1$.  Then $e^x=1$, so $x=0$.  \\
    The point is $(0, 1-e^0)=(0,0)$ so the tangent line is $y=-1(x)$}
    \begin{center}
        \desmos{dwxditai89}{800}{600}
    \end{center}
    \textcolor{blue}{Use the colored dots next to the formulas to toggle the graphs on/off.}
    \textcolor{orange}{Both this and the next problem require the simplest version of solving an exponential.  Of course, students can use logs here, but do not need to if they can see the solution by inspection.}
\end{exercise} 

\begin{exercise}
    At what value of $x$ does $f(x)=e^x-ex$ have a horizontal tangent line? \\
    \textcolor{blue}{ $f'(x)=e^x-e$\\
    A horizontal tangent has slope 0, so we need $e^x-e=0$.\\
    $e^x=e$ means $x=1$.}
\end{exercise}

\begin{exercise}
    Calculate the derivative for the following functions using rules from our class: 
        \begin{enumerate}
            \item $f(x)=(3x^7-\sqrt{x})e^x$\\
            \textcolor{blue}{ $f'(x)=(3x^7-\sqrt{x})e^x+e^x(21x^6-\frac{1}{2}x^{-\frac{1}{2}})$}
       
            \item $h(t)=e^t(3t-2)t^{\frac{3}{4}}$\\
            \textcolor{blue}{ $h'(t)=\frac{t^{\frac{3}{4}}(3e^t+e^t(3t-2))-e^t(3t-2)(\frac{3}{4}t^{-\frac{1}{4}})}{t^{\frac{3}{2}}}$}
       
            \item $f(t)=\frac{te^{2t}}{e^{3t}}$\\
            \textcolor{blue}{ $=\frac{t}{e^t};\;\; f'(t)=\frac{e^t-te^t}{e^{2t}}$}
        
            \item $p(x)=\frac{\sqrt{x}e^x}{x-3}$\\
            \textcolor{blue}{ $p'(x)=\frac{(x-3)(\sqrt{x}e^x+e^x(\frac{1}{2}x^{-\frac{1}{2}}))-\sqrt{x}e^x}{(x-3)^2}$}
        \end{enumerate}
\end{exercise}

\begin{exercise}
    The intensity (in candles) of light seen by the eye when a camera flashes at time t (in milliseconds) is modelled by the equation
        $$I(t)=\frac{100t}{e^t}.$$
      
    Go graph this on Desmos and explore the graph a bit.
    %\begin{center}
        \[
        \graph{100x(e^{-x})}
        \]
    %\end{center}
    %\xmalt A graph of the function $I(t)=\frac{100t}{e^t}$.
    Calculate 
    \begin{enumerate}
        \item[a.] $I'(0.5)$ 
        \textcolor{blue}{$=30.327$}
        \item[b.] $I'(2)$
        \textcolor{blue}{ $=-13. 534$}
        \item[c.] $I'(5)$
        \textcolor{blue}{ $=-2.695$}
        \item[d.] Explain what these values represent. 
        \textcolor{blue}{ These values are the the instantaneous rate of change in light intensity at the given millisecond. Thinking about the physical situation, do these values make sense?}
\end{enumerate}

\end{exercise}


\end{document}